\properlane{whole-person-and-desire}{The Gospel reforms the whole person and teaches desire}

\subsection{Why apparently slight teaching can change a life}

Jerome explains the mustard seed by comparing the first appearance of Gospel preaching with its eventual effect. The preaching of a crucified God can seem implausible beside impressive philosophical discourse. Its beginning may appear small to the person who judges teaching by eloquence or prestige. Received in a believing soul, however, it grows into an understanding through which that person lives. Jerome explicitly allows the image to concern both a single soul and the spread of the Gospel throughout the world.\footnote{Jerome, \textit{Commentary on Matthew}, II, on Matthew 13:31--32, PL 26, cols.~89--91.}

His account of the sower also permits attention to the hearer's activity. Some identify the man sowing with the Savior; others understand the individual receiving and tending the seed in his own heart. Jerome develops this second possibility in terms of the mind receiving the preaching and nourishing it with faith. The hearer does more than notice an attractive proposition. He allows the received teaching to become something understood, retained and fruitful. This gives the Collect's continued meditation a concrete relation to the Gospel: attention can be the human activity through which the seed is tended.

The expanded tree is an image of the abundance that follows. Jerome interprets its branches as the teachings in which believing souls, or powers devoted to God's service, find rest. These identifications extend the image beyond physical growth. They make understanding and contemplation part of the kingdom's fruit within a person. Yet the starting point remains the preaching received. The individual does not create a private kingdom out of his own ideas; faith tends a seed given through the Gospel.

Jerome's interpretation is consequently compatible with outward witness, but it asks a different initial question from Chrysostom's. Chrysostom asks how a few weak disciples can be effective among many. Jerome asks how teaching that looks small can become a source of life in the person who receives it. The distinction becomes sharper when they explain the three measures of meal.

\subsection{Three measures and the reform of the faculties}

For Jerome, the woman can signify apostolic preaching or the Church gathered from different nations. She brings the leaven of scriptural knowledge and understanding into the meal. He offers spirit, soul and body as one interpretation of the three measures. Their being leavened together signifies the agreement of the person, whose parts should no longer work against one another.\footnote{Jerome, \textit{Commentary on Matthew}, II, on Matthew 13:33, PL 26, col.~91.}

He then offers a second account of that agreement through reason, anger and desire. The Gospel enables reason to attain prudence, directs anger against vice, and turns desire toward virtue. This is an interpretation of how a human being is changed, rather than a requirement that the biblical image be reduced to a single anatomical scheme. The faculties have different tasks. Reason discerns the good; anger can resist what damages it; desire moves toward its possession. They are brought into agreement when each serves the good the Gospel teaches.

The consequence is more specific than a general exhortation to improve. Anger, for example, is given an object: vice. That directs moral attention toward what needs correction instead of granting hostility against a neighbor a religious title. Desire is also given an object: virtue. It is capable of being educated, so conversion involves learning what to seek as well as learning what to refuse. Prudence then concerns the judgment through which these movements become fitting action.

This makes the Collect's inclusion of words and deeds especially pertinent. Reflection can become detached from action when a person knows an obligation and remains governed by a contrary desire. Speech can also reveal that division: a person may speak respectfully about charity while treating someone with contempt. Read with Jerome's explanation, the Collect asks for the agreement needed to overcome such divisions. The prayer's request to God locates this integrity within grace, while its reference to speech and action provides a way to recognize its effects.

Chrysostom takes the three measures as a large quantity of meal; Jerome offers spiritual interpretations in which the Gospel brings the person's faculties into agreement. Chrysostom explains the disproportion between the small amount of leaven and the meal it affects; Jerome explains the completeness of the person's transformation.

Jerome also reports an interpretation connecting the three measures with belief in the Father, Son and Holy Spirit. He calls its sense pious but explicitly warns that uncertain explanations of parables cannot establish dogma.\footnote{Jerome, \textit{Commentary on Matthew}, II, on Matthew 13:33, PL 26, cols.~91--92.} That warning matters beside the Trinity Preface. The Preface confesses the distinction of persons and unity of essence from God's revelation. Three portions of flour do not constitute an account of three divine persons, and their mixing cannot explain divine unity. Purified understanding receives the Church's confession; the Gospel's household image is not a substitute for it.

\subsection{Understanding Scripture requires an answer from the hearer}

The Gospel ends by naming hidden things disclosed through parables. Jerome recognizes the quotation from Psalm 77(78), attributed to Asaph. That psalm recounts God's dealings with Israel and the people's responses. Its mysteries concern actual events through which God instructs his people. The use of a parable does not empty history of its reality. It asks the hearer to understand what is being disclosed through the account.\footnote{Jerome, \textit{Commentary on Matthew}, II, on Matthew 13:35, PL 26, cols.~92--93; Psalm 77(78).}

The preceding verse distinguishes the crowd hearing parables from the disciples who seek explanation. Jerome observes that the disciples question the Savior at home. Understanding therefore involves a continuing relation to the teacher. For a Christian reader, acquiring words without pursuing their meaning would leave the leaven of scriptural understanding largely unreceived. Inquiry is part of the obedience through which knowledge can change conduct.

The Epistle provides an outward test of this inward account. The Thessalonians have turned from idols to serve God. Their understanding has changed allegiance, and that allegiance has changed their conduct. They are known for faith, charity and endurance. Their interior conversion can be recognized by its relation to their neighbors and by the example it gives others. The prayer for whole-person reform concerns members of an assembly whose lives have consequences for other people.

\subsection{The cry that acknowledges the need for mercy}

Augustine explains why that conversion includes prayer from the depths. His exposition of Psalm 129(130) first distinguishes being in distress from knowing that one needs deliverance. The person who laughs at the summons to pray may appear prosperous while being more deeply trapped. Success has become a reason to dismiss God's concern with his conduct. The sinner who cries, by contrast, has begun to acknowledge both his situation and the possibility of mercy.\footnote{Augustine, \textit{Expositions of the Psalms}, 129.1--2.}

That distinction changes how the repeated petition at the Alleluia and Offertory can be understood. The cry does not require the person to have completed his reform before asking for help. It is itself part of turning toward God. The prayer names need in the presence of someone able to answer. It also resists an inward scrutiny that would end only in contemplation of one's failures: the voice is addressed to the Lord.

Bellarmine develops the same practical difficulty through two necessities for penitence. A person must recognize the misery of sin and know God's mercy. Without the first, he seeks no medicine. Without the second, he despairs of a cure. The psalm holds both before him: the question about who could stand under the exact reckoning of iniquities and the assurance that forgiveness is with God.\footnote{Bellarmine, \textit{Commentary on the Book of Psalms}, Psalm 129, explanations 3--5.} This gives precise content to a desire for change. Repentance includes sorrow, but its movement is toward the mercy that makes a different life possible.

The Introit's psalm verse can be heard in that relation. Augustine connects the turning of captivity with the forgiveness that follows in the full psalm. The captivity is understood spiritually as bondage to sin. Bellarmine places its full release at the resurrection, while acknowledging the partial liberty already received.\footnote{Augustine, \textit{Expositions of the Psalms}, 84.2--6; Bellarmine, \textit{Commentary on the Book of Psalms}, Psalm 84, explanation 1.} Thus someone can give thanks for liberation and still ask for renewal. The continued need for help does not make the earlier mercy unreal.

Jeremiah's promise of gathering also calls the people to seek God with the whole heart. Jerome's exposition connects that wholehearted search with prayer. The language of gathering can therefore be placed beside the desire for agreement within the person, while its literal subject remains the restoration promised to the exiles. Here the connection concerns the completeness of the response God asks: the person who seeks mercy cannot reserve a contrary purpose as exempt from God's address.

The full psalm from which the Gradual is taken describes suffering endured without abandoning God's covenant. Faithfulness can therefore coexist with lament and outward injury. Interior reform does not guarantee emotional ease or immunity from outward harm. The worshipper may praise God's help and still lament an injury; faithful people need not prove their progress by feeling untroubled. The need named in the repeated cry remains real.

\subsection{Sacrifice and the prospect of a changed life}

Augustine moves from the psalm's admission of sin to its reason for hope: propitiation is with God. He explains propitiation through sacrifice, identifying the sacrifice offered for us in Christ. The sinner's expectation of forgiveness therefore has an object beyond the intensity of his remorse. Christ has acted for those who cannot establish their innocence before God.\footnote{Augustine, \textit{Expositions of the Psalms}, 129.3.}

That explanation belongs to the full psalm, beyond the verses selected in the two chants. It illuminates their use at the offering because the Secret asks for the oblation's cleansing and renewing effects. Augustine does not comment on this Secret; Schuster does, identifying its four requests. Read together, the psalm exposition and the prayer relate hope of forgiveness to the sacrifice and ask that its effects extend into the worshippers' continued life.\footnote{Schuster, \textit{The Sacramentary}, I, p.~417.}

The second pair of petitions, governance and protection, makes the continued life especially clear. Cleansing does not leave a person without direction. Renewal is to be preserved in choices made after the prayer. Jerome's faculties provide a concrete account of what needs governance: judgments, anger and desire can either agree in the pursuit of virtue or pull against one another. The sacrifice's fruits are sought for the people who must make those choices.

The Offertory's repeated cry keeps this offering humble. Schuster explains that the gifts are themselves received from God. Such an account is consistent with Augustine's recognition that the sinner's hope rests on what Christ has offered. Neither the goods presented nor the prayer spoken makes the worshipper self-sufficient. The person's participation is real, and the source of the good remains God.\footnote{Schuster, \textit{The Sacramentary}, III, p.~188; Augustine, \textit{Expositions of the Psalms}, 129.3--5.}

\subsection{What a person learns to ask for}

The Communion's promise creates a difficult question for the person being formed in this way. Strong desire can be directed toward things that would injure the person or someone else. The assurance of being heard cannot make the mere strength of a desire its warrant. Mark's following command to forgive and Bede's explanation of faith through charity give the petition a moral setting. Those asking pardon are to forgive; those who lack extraordinary perfection can confidently ask for remission of sins and eternal life.\footnote{Mark 11:24--25; Bede, \textit{Commentary on Mark}, on 11:22--25, \textit{Et cum stabitis ad orandum}.}

This does not make forgiveness a technique for obtaining unrelated possessions. It brings the petitioner into the same relation of mercy that he asks God to extend toward him. A refusal to forgive is a contradiction within the prayer for pardon. In the present interpretation, Jerome's agreement of the faculties has this practical form: what the mouth asks of God and what the will intends toward the neighbor must be brought together.

The Postcommunion then describes desire after reception. Those who have been fed ask to continue wanting the gifts by which they live. Desire here is capable of improvement through receiving its proper object. The gift can deepen recognition of one's need for it. Schuster calls attention to the unusual pairing of nourishment and hunger in this prayer.\footnote{Schuster, \textit{The Sacramentary}, I, p.~417.} The person who receives heavenly food is asking for a disposition that will continue to turn him toward it.

Augustine's later exposition of Psalm 129 helps explain the final direction of that desire. He reads the morning watch through Christ's resurrection. Christ has risen in the flesh he received from humanity; the believer is to hope for what has already occurred in the Head. Forgiveness has been received, while bodily resurrection is still awaited. The two gifts are related without being confused.\footnote{Augustine, \textit{Expositions of the Psalms}, 129.4--8.}

He consequently distinguishes the hope of Israel from treating earthly abundance or bodily health as the highest good. Such goods can be sought from God, but the hope disclosed by the resurrection exceeds them. This distinction helps explain how prayer can remain confident without promising an identical temporal outcome to every petitioner. The Son awaited in the Epistle is already risen. The person's reformed desire is directed toward sharing his life, and the request after Communion concerns continued nourishment toward that end.

\begin{fourSenses}
\item[Literal.] Jesus describes leaven acting throughout the meal. The Collect joins reflection to speech and deeds; the Secret asks for cleansing, renewal, governance and protection; the Postcommunion asks the already nourished to continue desiring life-giving gifts.
\item[Allegorical.] In Jerome's interpretation, the Church communicates scriptural understanding that transforms those who receive it. In Augustine's exposition, Christ's sacrifice gives the sinner hope of forgiveness and his resurrection gives the body hope of life. Thanksgiving is offered to the one God confessed in Trinity.
\item[Moral.] Understanding is to become prudent action, anger opposition to vice and desire the pursuit of virtue. The cry for mercy includes honest recognition of sin, confidence in God's forgiveness and the willingness to forgive others.
\item[Anagogical.] The person's transformation awaits complete redemption, including the resurrection of the body. The hunger expressed after Communion is directed toward the life already manifested in the risen Son, whom the community awaits.
\end{fourSenses}
